\section{Rangkuman}
\label{sec:rangkuman-bab12}

Pada bab ini telah dibahas:

\begin{itemize}
    \item Pengambilan keputusan berjalan dalam dua langkah: satu
          instruksi menyiapkan \eng{flag}, lalu satu instruksi
          lompatan bersyarat membacanya. Lima \eng{flag} terpenting
          adalah \flag{ZF}, \flag{CF}, \flag{SF}, \flag{OF}, dan
          \flag{PF} (Bab~\ref{chap:register}).

    \item \asm{CMP} menyiapkan \eng{flag} melalui pengurangan yang
          hasilnya dibuang, sedangkan \asm{TEST} melalui \asm{AND}
          yang hasilnya dibuang. \asm{CMP} membandingkan nilai;
          \asm{TEST} memeriksa bit.

    \item \asm{JMP} memindahkan kendali tanpa syarat. Jarak lompatan
          menentukan ukuran instruksi: \eng{short} dua byte
          ($\pm$ 127 byte), \eng{near} tiga byte (satu segmen), dan
          \eng{far} lima byte (antarsegmen).

    \item Keluarga \asm{Jcc} terbagi menjadi kelompok bertanda
          (\asm{JG}, \asm{JGE}, \asm{JL}, \asm{JLE}), kelompok tak
          bertanda (\asm{JA}, \asm{JAE}, \asm{JB}, \asm{JBE}),
          pemeriksa \eng{flag} tunggal (\asm{JC}, \asm{JS},
          \asm{JO}, \asm{JP}, \asm{JCXZ}), dan kelompok kesamaan
          (\asm{JE}/\asm{JNE}) yang berlaku untuk keduanya.

    \item Perulangan dibentuk dengan \asm{LOOP} dan \reg{CX}, varian
          \asm{LOOPE}/\asm{LOOPNE} yang menambahkan syarat \flag{ZF},
          atau pola \asm{DEC} dan \asm{JNZ} yang lebih luwes dan
          lebih cepat pada prosesor modern.

    \item Seluruh struktur kontrol bahasa tingkat tinggi---\code{if},
          \code{while}, \code{for}, dan \code{switch}---diterjemahkan
          dari tiga bahan dasar: \asm{CMP}/\asm{TEST}, \asm{Jcc}, dan
          \asm{JMP}. \code{switch} dengan banyak cabang paling
          efisien memakai tabel lompatan.

    \item Studi kasus pencarian maksimum, pengurutan gelembung, dan
          pencarian biner menegaskan bahwa struktur kontrol kompleks
          tetap dibentuk dari bahan dasar yang sama.
\end{itemize}

\begin{checklistbox}
Tandai setiap butir yang sudah Anda kuasai.

\begin{itemize}
    \item[\checkbox] Saya dapat menjelaskan peran \eng{flag} dan
          menyebutkan instruksi yang menyiapkannya.
    \item[\checkbox] Saya dapat membedakan \asm{CMP} dan \asm{TEST}
          beserta situasi pemakaian masing-masing.
    \item[\checkbox] Saya dapat menjelaskan tiga jenis jarak lompatan
          dan pengaruhnya pada ukuran kode.
    \item[\checkbox] Saya dapat memilih \asm{Jcc} yang benar untuk
          data bertanda dan tak bertanda tanpa melihat tabel.
    \item[\checkbox] Saya dapat menulis perulangan dengan \asm{LOOP}
          maupun dengan \asm{DEC} dan \asm{JNZ}.
    \item[\checkbox] Saya dapat menerjemahkan pseudocode
          \code{if--else} dan \code{while} ke assembly.
    \item[\checkbox] Saya dapat menjelaskan cara kerja tabel lompatan
          pada pola \code{switch}.
    \item[\checkbox] Saya dapat menelusuri alur percabangan dengan
          \eng{Single Step} dan menjelaskan keputusan setiap
          lompatan.
\end{itemize}

\vspace{0.4em}
Bila masih ada butir yang belum bertanda, ulangi bagian yang
bersesuaian sebelum melanjutkan ke bab berikutnya.
\end{checklistbox}

\begin{infobox}
\textbf{Rujukan bab.} Materi pada bab ini diperkuat oleh
\cite{bartlett2004}, \cite{hyde2010}, \cite{evans2020}, dan
\cite{irvine2011}. Untuk penelusuran lebih dalam mengenai perulangan
dan pemilihan instruksi lompatan pada program nyata, lihat
\cite{irvine2011} dan \cite{hyde2010}. Penerapan tabel lompatan pada
program besar dibahas pada \cite{evans2020}.
\end{infobox}
